\section{系统功能划分与协作边界}\label{sec:division}
本报告以一台 TonyPi 和四台 TurtleBot2 组成的异构系统为对象。本组负责 TonyPi 抱物路线与方式二建图避障跟随。其余方法由对应小组独立补充，本稿保留技术接口与现有实现摘要，组内主要分工按成员提供的信息归纳，交叉参与的工作以协作方式说明。
\begin{table}[H]\centering\small
\begin{tabularx}{\linewidth}{p{27mm}YY}\toprule
功能模块 & 技术范围 & 协作输入与输出\\\midrule
TonyPi & 抱物动作、固定标记识别、路线状态机、转向与放置 & 场地 A/B/C 标记；动作、日志、目标外观与任务状态适配\\
方式二 & 建图、定位、AR目标转换、轨迹跟随、规划与门控 & 实测地图、定位、扫描与任务许可；安全底盘速度\\
方式一 & 当前提交的 AR 直接跟随实现及其验证 & 背部 ALVAR 标记、目标位姿与速度控制\\
方式三 & RGB-D 候选、身份关联及实时避障 & RGB-D 与本车里程计；目标、代价地图与控制\\
方式四 & 前车位姿或轨迹队列跟随 & 唯一前车身份、共用坐标与队列间距\\\bottomrule
\end{tabularx}\caption{系统功能与模块边界。}
\end{table}
各组在独立章节和图片子目录中修改技术内容。公共导言、主文件及引用格式统一维护；跨模块参数变化应同步记录接口影响。队列拓扑、跟随目标来源和导航方法分别登记，避免将同一控制器在不同车号上的运行视为不同方法。

\subsection{第一小组组内分工与协作}
第一小组采用共同设计、交叉调试与协同完善的工作方式，围绕动作标定、二维码转弯、状态机实现、实验环境和报告整理形成相互衔接的任务分工。表中归纳主要承担方向；共同参与的内容不重复计为互相独立的成果，也不作定量贡献比例分配。
\begin{longtable}{p{18mm}p{37mm}p{88mm}}
\caption{第一小组成员的主要分工与交叉协作}\\\toprule
成员 & 主要承担方向 & 工作内容与协作关系\\\midrule\endfirsthead
\toprule 成员 & 主要承担方向 & 工作内容与协作关系\\\midrule\endhead
\textbf{杨开元} & 动作标定、状态机与转弯算法初步实现 & 面向抱物行走与左右转弯开展动作标定，编写路线状态机并初步实现二维码转弯算法；衔接标记观测、阶段判断与动作调用，结合现场表现与成员共同完善转换条件和执行逻辑。\\
\textbf{林楗棋} & 动作组改进与转弯算法优化 & 调整与优化动作组及动作衔接，参与二维码转弯算法的实现与迭代；配合观察视觉判断与实际转向的对应关系，协同改进动作配置及算法细节，推进感知决策与机械执行的协调。\\
\textbf{曹淇涵} & 赛道搭建与动作辅助标定 & 参与围挡通道、转角及运行区域搭建和调整，协助开展动作标定与现场试验；配合核对机器人姿态、行走方向和转弯表现，将场地约束及运行现象反馈给标定与算法调试。\\
\textbf{刘小凡} & 算法协同构思、现场配合与报告撰写 & 参与二维码识别与转弯决策方案构思，协同开展动作标定、赛道搭建与现场调整；承担报告组织、技术机制梳理、图表制作和视频证据整理，将实现过程汇总为可复核的技术说明。\\\bottomrule
\end{longtable}
上述分工贯穿“方案讨论—动作标定—算法实现—场地验证—结果整理”全过程。动作和算法之间并非单向交付：转弯幅度影响视觉决策门限，围挡布局影响观测条件，现场反馈又推动动作组与程序调整。成员的主要方向彼此衔接，最终结果来自多轮共同讨论与交叉调试。
